% ---
% §7.4 Viabilidad práctica (OE-6).
% Entradas que lo respaldan: D-040, D-042, H-019, H-032, H-051, H-052, H-054, H-055.
% ---

\subsection{Viabilidad práctica}
\label{sec:conclusion_viabilidad}

% fuente: H-051, H-054; 8_oe6_qubits.csv, X8_costo_gradiente.csv
En simulación, el número de qubits no fue la restricción que condicionó el diseño. En la máquina
empleada, con 16\,GiB de memoria, el \textit{embedding} de todas las lesiones cabe en 24\,h hasta unos
25 qubits y el kernel de $200 \times 200$ por \textit{compute-uncompute} hasta unos 22. El límite lo
impuso el gradiente: entrenar los parámetros variacionales por \textit{parameter-shift}, la regla que
funciona en \textit{hardware}, cuesta 5.16\,s por lesión a $k = 12$, unas 205\,h para 50 épocas,
y deja de caber en un día a partir de unos 15 qubits.

% fuente: H-032, H-052, H-054, H-055; 8_oe6_shots_kernel.csv, 8_oe6_shots_embedding.csv, X9_costo_hardware.csv
% respaldo: «Thus on training with a polynomial number of measurements, one ends up with a trivial model where the predictions on unseen inputs are independent of the input data.» (Thanasilp2024, p. 1)
Varios de los atajos que hicieron viable el estudio solo existen en simulación. El kernel de
$200 \times 200$ se calculó en 1.5\,s a partir de los vectores de estado, y el gradiente por
retropropagación costó unas 3\,600 veces menos que por \textit{parameter-shift}; ninguno de los dos es
realizable en un procesador cuántico, porque medir colapsa el estado. En \textit{hardware}, cada entrada
del kernel es un circuito distinto que se repite con múltiples disparos: 19\,900 circuitos por matriz
con 200 lesiones y más de un millón con todas las calcificaciones de entrenamiento. Con 1\,024
disparos, el perceptrón sobre los $\langle Z_i\rangle$ queda a menos de 0.02 del AUC exacto, pero el
kernel no es medible: a $c = 1$, entre el 16 y el 74\,\% de sus valores leería exactamente cero. Es la
situación que describen Thanasilp et al. \cite{Thanasilp2024}, en la que con un número polinomial de
mediciones las predicciones dejan de depender de los datos. La viabilidad de estas codificaciones en un
contexto de salud queda condicionada, antes que al número de qubits, a evitar la concentración del kernel
y a disponer de gradientes que un procesador real pueda estimar.
